将 JHQ 移植到 GPU 会遇到三个技术挑战。

第一，压缩编码减少了存储，但扫描编码仍可能代价较高。每个八位主码索引一张 256 项的距离查找表（lookup table，LUT）；这些表需要构造，且可能无法放入快速的片上内存。由于 JHQ 通过组合各坐标的量化等级形成编码，这张表可以在实数运算下精确拆成两张 16 项表，分别对应一个字节的前、后四位。挑战在于：更小的表能否抵消额外查表和加法的开销。

第二，小批量可能只有少量查询，不足以充分利用 GPU。将每条查询的候选分给多个线程块可以增加并行工作，但如果各块仅保留相同份额，最好的候选又集中在一个块中，就可能丢失全局优胜者。一种办法是在归并前，让每组各自保留完整的候选短名单，代价是增加显存和归并开销。挑战在于保留所有可能胜出的候选，同时让这些开销低于并行带来的收益。本文研究共享截断阈值和随批量变化的调度，以降低这些开销。

第三，达到目标召回率所需的精排工作量随工作负载变化。候选过少可能丢失真实近邻，过多则可能花时间读取和计算残差编码，却不改善答案。目标是在不使用真值标签的情况下选择可复用预算。较大预算搜索可以提供参考，但它自身也可能漏掉真实近邻。因此，既要在独立测试查询上核查召回率，也要测量校准和后续搜索的成本。
